\subsection{The original adic coefficient map into actual compatible sections}

Retain ANY commutative ring $R$, ANY ideal $I\subseteq R$, ANY scheme
$X$, and ANY original morphism $f:X\to\operatorname{Spec}R$. Write
$X_n=X\times_{\operatorname{Spec}R}\operatorname{Spec}(R/I^n)$ with the
LITERAL original projections, and retain the actual transitions
$j_{n,m}:X_n\to X_m$ for $n\le m$. The actual compatible-section ring is
$C=\{(s_k)_{k\ge0}:s_k\in\Gamma(X_{k+1},\top),
 j_{k+1,l+1}^*(s_l)=s_k\text{ for every }k\le l\}$.
These are the already constructed whole schemes and section RingHoms.
No properness, Noetherianity, separation, reducedness, finiteness or
completion comparison hypothesis enters either producer below.

\begin{theorem}[Every original quotient coefficient map retains the whole transition square]
\label{scheme:actual-original-infinitesimal-coefficient-map}
\leanok
\lean{EnvelopeProperty.AdicCompletionSections.coefficient,
EnvelopeProperty.AdicCompletionSections.transition_coefficient,
EnvelopeProperty.AdicCompletionSections.restriction_coefficient,
EnvelopeProperty.AdicCompletionSections.coefficientFamily,
EnvelopeProperty.AdicCompletionSections.coefficientFamily_apply}
\uses{scheme:actual-original-infinitesimal-idempotent-family,scheme:actual-proper-global-section-valuation-membership}
Let $c_n:R/I^n\to\Gamma(X_n,\top)$ be the ORIGINAL coefficient RingHom,
obtained from the actual second projection by the genuine
$\Gamma(\operatorname{Spec}(R/I^n))\simeq R/I^n$ isomorphism. For EVERY
$n\le m$ prove the ENTIRE RingHom identity
\[
 j_{n,m}^*\circ c_m=c_n\circ\pi_{n,m},
 \qquad \pi_{n,m}:R/I^m\to R/I^n.
\]
The whole coefficient RingHom $c:R\to C$ is thereby constructed, with
EVERY actual coordinate exactly
$c(r)_k=c_{k+1}(r\bmod I^{k+1})$. Its coordinate maps and original
$R$-coefficient laws are proved on every element, not just on supports.
\end{theorem}
\begin{proof}
Use the already constructed WHOLE scheme identity
$\operatorname{pr}_{2,m}\circ j_{n,m}
 =\operatorname{Spec}(\pi_{n,m})\circ\operatorname{pr}_{2,n}$.
Apply the genuine contravariant global-section functor. The ACTUAL
naturality of the original $\Gamma\operatorname{Spec}$ isomorphisms
identifies the coefficient corners with the SAME quotient rings. This
gives the displayed RingHom identity on every original quotient element.
Compose it with the ORIGINAL quotient maps $R\to R/I^m$; their actual
factor law gives every compatibility equation for the tuple defining
$c(r)$. Each coordinate is an actual RingHom, so coordinatewise ring
operations construct the true RingHom $c$ into the ORIGINAL subring $C$.
\end{proof}

\begin{theorem}[The original adic completion has its actual coefficient RingHom into the whole section limit]
\label{scheme:actual-adic-completion-compatible-section-coefficients}
\leanok
\lean{EnvelopeProperty.AdicCompletionEvalCompatibility.factorPow_eval,
EnvelopeProperty.AdicCompletionEvalCompatibility.factorPow_comp_eval,
EnvelopeProperty.AdicCompletionSections.completionCoefficient,
EnvelopeProperty.AdicCompletionSections.completionCoefficient_apply,
EnvelopeProperty.AdicCompletionSections.completionCoefficient_of,
EnvelopeProperty.AdicCompletionSections.completionCoefficient_comp_algebraMap,
EnvelopeProperty.AdicCompletionSections.idempotentFamily,
EnvelopeProperty.AdicCompletionSections.idempotentFamily_zero}
\uses{scheme:actual-original-infinitesimal-coefficient-map,scheme:actual-original-infinitesimal-idempotent-family}
Let $\widehat R=\operatorname{AdicCompletion}_I(R)$ be Mathlib's ACTUAL
adic completion. Construct the ORIGINAL coefficient RingHom
$\widehat c:\widehat R\to C$ with EVERY genuine coordinate formula
\[
 \widehat c(a)_k=c_{k+1}(\operatorname{ev}_{k+1}(a)),
 \qquad \operatorname{ev}_n:\widehat R\to R/I^n.
\]
Prove the WHOLE RingHom equality
$\widehat c\circ\iota_R=c$, where $\iota_R:R\to\widehat R$ is the actual
original completion coefficient map. Actual coordinate evaluation and
EVERY original section restriction commute with $\widehat c$ through the
SAME original quotient maps. A true idempotent of $\widehat R$ maps to
an actual idempotent compatible family, whose original order-one
coordinate is exactly $c_1(\operatorname{ev}_1(a))$.
No injectivity, surjectivity, image characterization or isomorphism of
$\widehat c$ is asserted or assumed.
\end{theorem}
\begin{proof}
REUSE the genuine existing algebra RingHoms
$\operatorname{ev}_n$ of the original adic completion; do not define a
replacement completion. Their actual compatible-tuple property, with
the genuine equality $I^n\top=I^n$ transported through the original
quotient RingEquivs, gives
$\pi_{n,m}(\operatorname{ev}_m(a))=\operatorname{ev}_n(a)$ for EVERY
$n\le m$. Combine this WHOLE quotient identity with the preceding
original coefficient transition square. The resulting tuple therefore
belongs to the ACTUAL compatible-section subring $C$. Every coordinate
is the actual composite of two RingHoms, so the whole tuple map is a
RingHom and has exactly the displayed coordinates. The existing genuine
identity $\operatorname{ev}_n(\iota_R(r))=r\bmod I^n$ gives every
coordinate of $\widehat c\circ\iota_R=c$; true subtype/function
extensionality gives the WHOLE map equality. RingHom preservation of
idempotents and the actual first coordinate yield the final laws.
\end{proof}

\begin{remark}[The exact formal-functions obligation remains]
This constructs a map INTO the section limit; it does not prove a
completion comparison. In the proper Noetherian-local situation the
source proof still needs the genuine formal-functions theorem identifying
the completed pushforward stalk with this ACTUAL section limit, together
with the finite-module completion comparison, before any arbitrary fibre
idempotent can be shown to lie in the original pushforward-stalk image.
Existing localness of the actual finitely generated maximal-adic
completion is REUSED when needed, not postulated or reproved here.
No proper-fibre idempotent image theorem, Stein factorization,
geometrically connected proper fibre, or native connectedness conclusion
follows from the two coefficient-map producers alone.
\end{remark}
